On-policy distillation (OPD) improves student models by training them on trajectories induced by their own policy, making it a promising approach for mitigating exposure bias in agent training.
%
However, most OPD studies focus on single-turn settings, while realistic LLM agents interact with environments over multiple turns.
%
% \zz{
In this regime, early errors can alter future observations and compound across the trajectory, and standard dense token-level OPD becomes brittle, as it may over-penalize semantically valid alternatives, reinforce local degeneracies such as repeated actions, and propagate unreliable teacher supervision on off-distribution histories.
% }
%
We propose \textbf{\ours}, a verifier-free selective intervention framework specifically designed for multi-turn OPD.
%
Instead of applying teacher supervision uniformly across all turns, \ours first observes environment feedback and uses teacher judgment to decide whether each student response should be skipped or intervened on.
%
To further address compounding errors, \ours weights token-level distillation by teacher confidence, reducing the influence of uncertain teacher distributions on corrupted or ambiguous histories.
%
Finally, \ours applies loss normalization to preserve the overall loss scale of standard OPD while retaining selective turn-level weighting.
%
Experiments on agent tasks show that \ours consistently improves over baselines, achieving up to a $13.3\%$ relative improvement in ALFWorld unseen success rate over standard OPD.
%
Ablation studies further demonstrate that turn-level intervention, teacher confidence weighting, and loss normalization provide complementary benefits.
%
Our results suggest that effective multi-turn OPD should remain on-policy, but teacher supervision should be selectively allocated to turns where intervention is necessary and reliable.
%
